\documentclass[10pt]{article}
\usepackage[margin=1in, paperwidth=8.5in, paperheight=11in]{geometry}
\usepackage{ifpdf,amsmath, amssymb, comment, color, graphicx, stmaryrd,setspace,enumitem,tikz, fancyhdr, wrapfig, textcomp, units, mathptmx, pdfpages, cancel, multicol}
\usepackage{hyperref}
\hypersetup{
    colorlinks=true,
    urlcolor=blue,
}

\setlength{\headheight}{14.5pt}
\newcommand{\Q}{\mathbb{Q}}
\newcommand{\R}{\mathbb{R}}
\newcommand{\Z}{\mathbb{Z}}
\newcommand{\vu}{\mathbf{u}}
\newcommand{\vv}{\mathbf{v}}
\newcommand{\vw}{\mathbf{w}}
\newcommand{\vi}{\mathbf{i}}
\newcommand{\vj}{\mathbf{j}}
\newcommand{\vk}{\mathbf{k}}
\newcommand{\vn}{\mathbf{n}}
\newcommand{\vr}{\mathbf{r}}
\newcommand{\va}{\mathbf{a}}
\newcommand{\vF}{\mathbf{F}}
\newcommand{\vL}{\mathbf{L}}
\newcommand{\vT}{\mathbf{T}}
\newcommand{\vN}{\mathbf{N}}
\newcommand{\proj}{\operatorname{proj}}
\newcommand{\orth}{\operatorname{orth}}
\newcommand\del\nabla
\newcommand\dotp[1][.5]{\,\mathbin{\vcenter{\hbox{\scalebox{#1}{$\bullet$}}}}\,}
\newcommand{\slice}{_\text{slice}}

% Solution text is in red. If you want the solutions to show, remove the \iffalse from the definition of the \red command.
%\newcommand{\red}[1]{ %\iffalse
%	\textcolor{red}{#1} }%\fi}

\newenvironment{red}{\color{red}}{\ignorespacesafterend}
\newenvironment{blue}{\color{blue}}{\ignorespacesafterend}
\newcommand{\green}[1]{\textcolor{green}{#1}}
\renewcommand{\section}[1]{\begin{center} \textbf{#1} \\\end{center}}
%
\hyphenpenalty=5000
\setlength{\parindent}{0in}
%\oddsidemargin=-.25in
\allowdisplaybreaks
\pagestyle{fancy}
\renewcommand{\headrulewidth}{0pt}
\lhead{MATH 202}
\rhead{Spring 2024}
%\lfoot{\copyright\ CLEAR Calculus 2010}
\cfoot{}

\begin{document}
%


%\onehalfspacing
\allowdisplaybreaks
%##################################################################
\section{PS\#8: Improper integrals - \red{Answer key} }

\everymath{\displaystyle}
\begin{enumerate}[leftmargin=0pt]

\item (\href{https://activecalculus.org/single/sec-6-5-improper.html#IZN}{Activity 6.5.4}) For each of the following definite integrals, decide whether the integral is improper or not. If the integral is proper, evaluate it using the First FTC. If the integral is improper, determine whether or not the integral converges or diverges; if the integral converges, find its exact value.
\begin{enumerate}
    \item $\displaystyle \int_0^1 \frac{1}{x^{1/3}} \, dx$
    
    \begin{red}
    This integral is \textbf{improper} because it is infinitely tall: the integrand has a vertical asymptote at $x=0$, which is included in the bounds of integration. Since this is a $p$-integral with $p < 1$, I know the integral is going to \textbf{converge}. Let's figure out what it converges to:
    %
    \begin{align*}
        \int_0^1 \frac{1}{x^{1/3}} &= 
        \lim_{a\to 0}\int_a^1 \frac1{x^{1/3}}\, dx =
        \lim_{a\to 0}\int_a^1 x^{-1/3}\, dx \\
        &= \lim_{a\to 0} \left[ \frac{x^{2/3}}{2/3} \right|_a^1 
        = \lim_{a\to 0}\left[\frac{1^{2/3}}{2/3} - \frac{a^{2/3}}{2/3}\right] = \frac{3}{2} - 0 = \frac{3}{2}.
    \end{align*}
    \end{red}
    \item $\displaystyle \int_0^2 e^{-x} \, dx$
    
    \begin{red}
    This integral is \textbf{proper}: it is neither infinitely long nor infinitely tall. So it definitely \textbf{converges}, and finding its value is straightforward:
    \[\int_0^2 e^{-x} \, dx = \left.\vphantom{\int} -e^{-x}\right|_0^2 = -e^{-2} - (-e^{-0}) = -e^{-2}+1.\]
    \end{red}
    \item $\displaystyle \int_1^4 \frac{1}{\sqrt{4-x}} \, dx$
    
    \begin{red}
    This integral is \textbf{improper} because it is infinitely tall: the integrand has a vertical asymptote at $x=4$, which is included in the bounds of integration. I'm pretty sure this is going to \textbf{converge} because the power is $1/2 < 1$, but in order to figure out what it converges to, we'll need to do a substitution $u = 4-x$, $du = -dx$:
    \begin{align*}
        \int_1^4 \frac{1}{\sqrt{4-x}} \, dx &= \lim_{b\to 4}\int_1^b \frac{1}{\sqrt{4-x}}\, dx \\
        &= \lim_{b\to 4}\int_{x=1}^{x=b} \frac{1}{\sqrt{u}}\, (-du) 
        = -\lim_{b\to 4}\int_{x=1}^{x=b} u^{-1/2}\, du \\
        &= -\lim_{b\to 4}\left[ \frac{u^{1/2}}{1/2} \right|_{x=1}^{x=b}
        \intertext{Now we back-substitute before plugging in the bounds of integration so that we do not make a puppy cry:}
        &= -\lim_{b\to 4}\left[ \frac{(4-x)^{1/2}}{1/2} \right|_1^b 
        = -\lim_{b\to 4}\left[ \frac{(4-b)^{1/2}}{1/2} - \frac{3^{1/2}}{1/2} \right]
        = -(0 - 2\sqrt{3}) \\
        &= 2\sqrt{3}.
    \end{align*}
    \end{red}
    
    \pagebreak
    
    \item $\displaystyle \int_{-2}^2 \frac{1}{x^2} \, dx$
    
    \begin{red}
    This integral is sneakily \textbf{improper} because the integrand has a vertical asymptote at 0, which is included in the bounds of integration, so the region is infinitely tall. Since this is a $p$-integral with $p > 2$, the integral will \textbf{diverge}. Let's see why:
    \[\int_{-2}^2 \frac{1}{x^2}\, dx = \int_{-2}^0 \frac{1}{x^2}\, dx + \int_0^2 \frac{1}{x^2}\, dx\]
    Here I'm going to play a trick: Look at a graph of $\frac{1}{x^2}$ and notice that it is symmetric across the $y$-axis. Therefore, these two integrals I've set up will have the same value, so I'm just going to calculate one of them and double it.
    \[ = 2\int_0^2 \frac{1}{x^2}\, dx = 2\lim_{a\to 0} \int_a^2 x^{-2}\, dx = 2\lim_{a\to 0} \left.\frac{x^{-1}}{-1} \right|_a^2 = 2\lim_{a \to 0}\left(\frac{-1}{2} - \frac{-1}{a}\right) \]
    In the limit as $a$ approaches 0, $\frac{1}{a}$ will increase without bound. Therefore, this integral \textbf{diverges}. 
    \end{red}
    \item $\displaystyle \int_0^{\pi/2} \tan(x) \, dx$
    
    \begin{red}
    This integral is sneakily \textbf{improper} because the integrand has a vertical asymptote at $\pi/2$, which is included in the bounds of integration, so the region is infinitely tall. If you happen to remember that $\int\tan(x)\, dx = \ln|\cos(x)| +C$, you could use that fact, but let's do it by substitution because that's more fun:
    %
    \begin{align*}
        \int_0^{\pi/2} \tan(x)\ dx 
        &= \int_0^{\pi/2} \frac{\sin(x)}{\cos(x)}\ dx 
        = \lim_{b\to \pi/2}\int_0^b\frac{\sin(x)}{\cos(x)}\ dx
        \intertext{So let's substitute $u = \cos(x)$, $du = \sin(x)\ dx$:}
        &= \lim_{b\to \pi/2}\int_{x=0}^{x=b} \frac{1}{u}\ du 
        = \left.\lim_{b\to \pi/2} \ln|u| \right|_{x=0}^{x=b}
        = \left.\lim_{b\to \pi/2} \ln|\cos(x)| \right|_0^b \\
        &= \lim_{b\to \pi/2} \left(\vphantom\int \ln|\cos(b)| - \ln|\cos(0)|\right)
        = \lim_{b\to \pi/2} \left(\vphantom\int\ln|\cos(b)| - \ln(1)\right) \\
        &= \lim_{b\to \pi/2} \left(\vphantom\int \ln|\cos(b)| - 0\right)
    \end{align*}
    And now we see a problem: as $b\to \pi/2$, $\cos(b) \to 0$, and so $\ln|\cos(b)| \to -\infty$. Therefore this integral \textbf{diverges}.
    \end{red}
    \item $\displaystyle \int_0^1 \frac{1}{\sqrt{1-x^2}} \, dx$
    
    \begin{red}
    This integral is \textbf{improper}. The region is infinitely tall because the integrand has a vertical asymptote at $x=1$, which is included in the bounds of integration. 
    %
    \begin{align*}
        \int_0^1\frac{1}{\sqrt{1-x^2}}\ dx 
        &= \lim_{b\to 1} \int_0^b \frac{1}{\sqrt{1-x^2}}\ dx = \lim_{b\to 1}\left[\vphantom\int \arcsin(x)\right|_0^b \\
        &= \lim_{b\to 1} \left[\vphantom\int \arcsin(b) - \arcsin(0) \right] 
        = \arcsin(1) - \arcsin(0) = \frac{\pi}{2} - 0 = \frac{\pi}{2}.
    \end{align*}
    Therefore this integral \textbf{converges}.
    \end{red}
\end{enumerate}

\end{enumerate}
\end{document}

\pagebreak

\item (\href{https://activecalculus.org/single/sec-6-5-improper.html#Hfn}{AC 6.5 Exercise 7})
Sometimes we may encounter an improper integral for which we cannot easily evaluate the limit of the corresponding proper integrals. In such cases, we often compare the difficult integrand to a simpler integrand. We use this idea in the exercises that follow.

\begin{enumerate}
    \item Explain why $x^2 + x + 1 > x^2$ for all $x \geq 1$, and hence show that $\int_1^\infty \frac{1}{x^2 + x + 1}\ dx$ converges by comparison to $\int_1^\infty \frac1{x^2}\ dx$.
    
    \begin{red}
    Since adding a (positive) number to $x^2$ produces something bigger than $x^2$, and since $x+1$ is certainly positive whenever $x \geq 1$, $x^2 + x + 1 > x^2$.
    
    Therefore, $\dfrac{1}{x^2 + x + 1} < \dfrac{1}{x^2}$, because a larger denominator produces a smaller fraction. 
    
    Therefore, $\displaystyle \int_1^\infty \frac{1}{x^2 + x + 1}\ dx < \int_1^\infty \frac{1}{x^2}\ dx$. This second integral is convergent because it's a $p$-integral with $p>1$.
    
    Thus, since the first integral is smaller than a convergent integral, there's no way it can diverge -- it must converge!
    \end{red}
    
    \item Observe that for each $x>1$, $\ln(x) < x$. Explain why \[\int_2^b \frac{1}{x} \, dx \ <  \int_2^b \frac{1}{\ln(x)} \,dx\] for each $b > 2$. Why must it be true that $\int_2^b \frac{1}{\ln(x)} \,dx$ diverges?
    
    \begin{blue}
    \textbf{NOTE:} There's a typo in this problem! The last line should say:\\
    ``Why must it be true that $\int_2^\infty \frac{1}{\ln(x)} \,dx$ diverges?'' (That top bound on the integral should be $\infty$ rather than $b$.)
    \end{blue}
    
    \begin{red}
    Since $x > \ln(x)$, $\frac{1}{x} < \frac{1}{\ln(x)}$, because the larger denominator produces the smaller fraction. 
    
    So let's think about the integrals $\int_2^b \frac{1}{x} \, dx $ and $  \int_2^b \frac{1}{\ln(x)} \,dx$. As long as $b > 2$, the interval we're integrating over is in this region $x > 1$ where $\ln(x) < x$. Therefore, $\int_2^b \frac{1}{x} \, dx \ <  \int_2^b \frac{1}{\ln(x)} \,dx$.
    
    Taking the limit as $b\to\infty$, we have $\int_2^\infty \frac{1}{x} \, dx \ <  \int_2^\infty \frac{1}{\ln(x)} \,dx.$
    
    We know that $\int_2^\infty \frac{1}{x} \, dx$ diverges because it's a $p$-integral with $p = 1$. Therefore, since $\int_2^\infty \frac{1}{\ln(x)} \,dx$ is \textbf{even bigger}, it has to diverge as well.
    \end{red}
    
    \item Explain why $\sqrt{\frac{x^4+1}{x^4}} > 1$ for all $x > 1$. Then, determine whether the improper integral $\int_1^{\infty} \frac{1}{x} \cdot \sqrt{\frac{x^4+1}{x^4}} \, dx$ converges or diverges.
    
    \begin{red}
    If we add 1 to something, we get something bigger, so $x^4+1 > x^4$. Dividing a number by something smaller than it produces something bigger than 1, so $\frac{x^4+1}{x^4} > 1$. Also, the square root of something that's bigger than 1 is for sure bigger than 1, so $\sqrt{\frac{x^4+1}{x^4}} > 1$.
    
    Multiplying a number by something bigger than 1 produces an even larger number. Therefore, $\frac1x < \frac1x \cdot\sqrt{\frac{x^4+1}{x^4}}$, so $\int_1^\infty \frac1x \ dx < \int_1^\infty \frac1x \cdot\sqrt{\frac{x^4+1}{x^4}}\ dx$.
    
    Similar to our reasoning in the last problem, the first integral diverges because it's a $p$-integral with $p=1$, so the even larger second integral also must diverge.
    \end{red}
\end{enumerate}

\end{enumerate}
\end{document}